\chapter{迁移学习}
\label{chap:transfer}
\begin{quote}\small
\textit{本章摘要。}迁移学习研究如何把源域中已经获得的数据、表示或参数用于目标域，以减少目标标签、训练时间或计算资源。本章先说明迁移学习的对象和历史分工，再区分任务适配、领域适配与领域泛化。随后从偏移假设出发，推导重要性加权、统计对齐和域对抗方法，最后讨论微调、Adapter、LoRA、测试时适配、负迁移及评价协议。
\end{quote}

本章主要参考 Jiang、Shu、Wang 和 Long 的综述\textit{Transferability in Deep Learning: A Survey}（2022），采用其中“预训练获取知识、适配利用知识、评价检验迁移”的区分。\footnote{Junguang Jiang, Yang Shu, Jianmin Wang, Mingsheng Long. \textit{Transferability in Deep Learning: A Survey}. 2022. arXiv:2201.05867。本文使用项目所附的 2022 年版本，\href{https://arxiv.org/abs/2201.05867}{公开版本入口}。}综述把统计匹配与域对抗作为深度领域适配的主要机制。本章选取其中适合初学者推导的方法，并补充偏移假设、低秩适配与部署评价。领域泛化是与适配有关但信息条件不同的设置，不作为所有迁移问题的共同前提。

第\ref{chap:pretraining}章说明预训练如何产生可复用的编码器。预训练参数只是一种可利用的资源，它不保证在新设备、新任务或新类别上仍然有效。本章回答一个更具体的问题：当数据产生环境或预测目标发生变化时，哪些源知识可以安全地使用，怎样使用，怎样判断迁移确实降低了目标风险？这一问题也是迁移学习综述所说的可迁移性（transferability）问题\cite{pan2010transfer}。

设旧相机拍摄的缺陷图像很多，新相机只有少量标注。旧模型可能已经学会边缘和纹理。新相机的亮度分布可能不同，缺陷类别也可能增加。如果把所有差异都称为“域偏移”，就无法决定应当调整样本权重、表示还是输出头。本章按“问题—对象—假设—机制—局限—评价”的顺序建立这一区分。第\ref{chap:continual}章将继续讨论目标数据还会按时间到达时，如何适配新情况并保持旧能力。

\section{迁移问题与知识来源}

\subsection{预训练与适配的分工}
迁移学习通常包含两个阶段。预训练（pre-training）在上游数据和上游任务上学习参数或表示。适配（adaptation）把这些参数或表示用于下游任务，必要时用目标数据更新它们。综述中的两阶段图景如图\ref{fig:transfer-lifecycle}所示。预训练追求对许多下游任务都有用的通用可迁移性。适配针对一个具体任务或一个具体目标域，追求在当前条件下的目标风险下降。两者可以由同一个训练过程连接，但训练目标和评价对象仍然不同。

\begin{figure}[htbp]
\centering
\begin{tikzpicture}[node distance=0.65cm and 0.65cm]
\node[idi input,text width=2.1cm] (up) {上游数据\\有标签或无标签};
\node[idi process,text width=2.1cm,right=of up] (pre) {预训练\\学习 $\theta_0$};
\node[idi state,text width=2.1cm,right=of pre] (model) {可复用模型\\$f_{\theta_0}$};
\node[idi input,text width=2.1cm,below=of model] (target) {目标数据\\有标签或无标签};
\node[idi process,text width=2.1cm,left=of target] (adapt) {适配\\更新部分参数};
\node[idi output,text width=2.1cm,left=of adapt] (eval) {目标测试\\$R_t$};
\draw[idi arrow] (up)--(pre)--(model);
\draw[idi arrow] (model)--(adapt);
\draw[idi arrow] (target)--(adapt);
\draw[idi arrow] (adapt)--(eval);
\end{tikzpicture}
\caption{预训练与适配的生命周期。箭头表示数据或模型参数进入下一阶段。目标测试标签只用于最终评价，不参与适配。}
\label{fig:transfer-lifecycle}
\end{figure}

预训练和迁移不是同义词。监督预训练、自监督预训练和领域继续预训练都描述参数如何先被学习。迁移学习描述这些参数如何被复用以及复用后是否改善目标任务。冻结编码器只训练新头、全量微调和低秩更新都是适配操作。第\ref{chap:pretraining}章中的编码器 $f_{\theta_0}$ 是本章的一个输入，而不是本章的结论。

\subsection{技术由来与迁移条件}
迁移学习的动因是逐个任务重新收集标签和训练模型的成本。早期统计方法已经研究训练样本与实际应用人群不同的问题。重要性加权通过改变样本贡献来修正采样偏差，核方法则在特征空间比较分布\cite{huang2007kmm}。这些方法说明，利用源知识可以从数据层面开始，不要求先有一个大型预训练网络。

深度表示学习使知识复用的位置发生了变化。监督视觉预训练把源类别识别中学到的特征用于检测等下游任务，自监督和语言建模又把预训练信号扩展到无标签数据。第\ref{chap:pretraining}章已经解释这些训练信号。本章关心训练结束之后的复用。源任务的输出头往往需要替换，骨干参数则可以冻结或微调。预训练与适配因而成为常见分工，但迁移学习在概念和方法上都早于这种深度训练流程\cite{pan2010transfer}。

这种分工留下两个不同困难。有目标标签时，需要控制目标拟合与可训练容量之间的取舍。没有目标标签时，仍然需要解释为什么源域的标签可以用于目标域。统计对齐和域对抗主要回应第二个困难\cite{long2015dan,ganin2016dann}。大型骨干又增加了保存和训练每个目标版本的代价，Adapter 与 LoRA 进一步限制目标专用参数\cite{houlsby2019adapter,hu2021lora}。这些机制可以组合，不能仅按模型年代替代对数据条件的检查。

\subsection{域、任务与分布}
设输入空间为 $\mathcal X$，标签空间为 $\mathcal Y$。一个任务规定输入要预测的对象、输出空间及损失函数。用 $f:\mathcal X\to\mathcal Y$ 表示理想的标注关系。一个域还包括数据产生环境和联合分布。源域记为 $\mathcal D_s$，目标域记为 $\mathcal D_t$，其联合分布分别为 $P_s(X,Y)$ 和 $P_t(X,Y)$。这里 $X$ 是随机输入，$Y$ 是随机标签。

任务适配（task adaptation）利用预训练模型和目标训练数据学习下游预测器。上游和下游可以有不同标签含义、标签空间或损失。例如从物体分类预训练转到缺陷严重度预测，需要新的回归头。任务适配也可以沿用原来的标签集合，名称不要求标签空间一定改变。综述中的这一设置主要提供预训练参数与带标签目标数据，不默认能重新访问全部上游数据。

领域适配（domain adaptation，DA）主要研究同一预测任务跨分布复用的问题，例如在不同相机或工厂上识别同样的缺陷。无监督领域适配（unsupervised domain adaptation，UDA）提供带标签源数据和无标签目标输入。半监督领域适配还提供少量目标标签。这些术语中的“无监督”或“半监督”描述目标域标签条件，源域通常仍有监督。

领域泛化（domain generalization，DG）要求模型在训练时没有见过的目标域上工作。它不能用目标输入估计密度比或目标统计量。常见训练条件是多个带标签源域，以及这些域之间存在可共享预测关系的假设。单源领域泛化还需要其他可解释的变化假设。对未见域的有效性必须通过留出整个环境的评价来检查。

这三者描述信息条件和目标，不是互斥的网络架构。一个模型可以先经过任务适配，再在目标无标签输入上做领域适配。若训练已经使用了目标域的无标签记录，该试验就不能同时声称检验了训练时完全未见该域的领域泛化。

图\ref{fig:transfer-source-target}以训练时的信息可见性比较三种设置。相同的网络和微调操作，在不同数据权限下属于不同实验。
\begin{figure}[htbp]
\centering
\begin{tikzpicture}[node distance=0.6cm and 0.6cm,font=\small]
\node[idi input,text width=3.5cm] (ta) {预训练参数 $\theta_0$\\目标标签与输入};
\node[idi process,text width=3.3cm,right=of ta] (tao) {任务适配\\学习目标预测器};
\node[idi input,text width=3.5cm,below=of ta] (da) {带标签源数据\\无标签目标输入};
\node[idi process,text width=3.3cm,right=of da] (dao) {无监督领域适配\\利用当前目标统计};
\node[idi input,text width=3.5cm,below=of da] (dg) {一个或多个源域\\目标输入训练时不可见};
\node[idi process,text width=3.3cm,right=of dg] (dgo) {领域泛化\\检验未见环境};
\draw[idi arrow] (ta)--(tao);
\draw[idi arrow] (da)--(dao);
\draw[idi arrow] (dg)--(dgo);
\end{tikzpicture}
\caption{三种迁移相关设置的典型训练信息。箭头表示可用信息进入学习过程。前两种可以使用相同的预训练骨干，无监督领域适配还可以同时利用源域标签和目标输入。领域泛化的目标域整体留出。}
\label{fig:transfer-source-target}
\end{figure}

同一设备不必构成一个任务。同一任务也可能由多个域产生。测试时是否提供设备型号、域标识或任务标识，会改变预测问题。若模型实际使用 $z$ 表示环境信息，预测规则应写为 $\hat y=f_\theta(x,z)$。若部署时没有 $z$，评价时也不能用真实域标识替代模型的输入。

\subsection{偏移假设}
迁移方法依赖关于 $P_s$ 和 $P_t$ 的假设。协变量偏移（covariate shift）假定
\begin{equation}
 P_s(Y\mid X)=P_t(Y\mid X),\qquad P_s(X)\ne P_t(X).
\label{eq:transfer-covariate-shift}
\end{equation}
输入出现频率变了，给定输入后的标签规则保持不变。标签偏移（label shift）假定
\begin{equation}
 P_s(X\mid Y)=P_t(X\mid Y),\qquad P_s(Y)\ne P_t(Y).
\label{eq:transfer-label-shift}
\end{equation}
类条件输入分布保持不变，类别比例发生变化。概念偏移（concept shift）表示 $P(Y\mid X)$ 改变。它可能来自标签定义变化、传感器与标签关系变化或未观测环境变量变化。

实际数据可能同时含有多种偏移。只有目标无标签输入时，通常不能仅凭数据证明式\eqref{eq:transfer-covariate-shift}成立。因此，偏移名称是建模假设，不能由文件夹名称直接确定。适配前应记录标签定义、采样条件、目标数据可见范围和标签到达时间。

\subsection{未见目标域的训练约束}
DG 没有目标输入，后文的目标密度比与目标统计量因而不可计算。一个可实施的起点是按源域组织风险。设有 $M\geq1$个带标签源域，第 $e$个源域的总体风险为 $R_e(f)=\mathbb E_{P_e}\ell(f(X),Y)$，$e=1,\ldots,M$。域均衡目标为
\begin{equation}
 \min_f\frac1M\sum_{e=1}^M R_e(f).
\end{equation}
实际用每个源域的训练平均损失替代 $R_e$。按域平均可避免大域仅凭样本量占据训练目标，但它只优化源域平均表现。若还希望控制最差源域，可以比较
\begin{equation}
 \min_f\max_{1\leq e\leq M}R_e(f).
\end{equation}
这里内部的最大值选择风险最高的源域，外部的最小化要求降低这一风险。后者会提高困难域的影响，也可能过度关注噪声较大的域。

这两个目标都不自动保证未知目标域表现。只有当目标变化与源域覆盖的变化有关，且存在可共享的预测关系时，多源训练才有可解释的迁移依据。源域留一验证可以模拟未见环境，却不能证明任意未来分布都被覆盖。采用保持标签的增强时，也必须解释增强是否符合真实设备变化。若增强删除了缺陷颜色或故障幅值，所谓不变性可能恰好损害目标任务。

\section{目标风险与重要性加权}

\subsection{从源风险得到目标风险}
给定预测函数 $f$ 和单样本损失 $\ell$，目标风险为
\begin{equation}
 R_t(f)=\mathbb E_{(X,Y)\sim P_t}\left[\ell(f(X),Y)\right].
\label{eq:transfer-target-risk}
\end{equation}
$R_t$ 是目标域上平均损失，不能由源测试分数直接替代。若满足协变量偏移且目标分布的输入支撑被源分布覆盖，则
\begin{align}
R_t(f)
&=\mathbb E_{P_s}\left[w(X)\ell(f(X),Y)\right],\\
 w(x)&=\frac{p_t(x)}{p_s(x)}.
\label{eq:transfer-importance-weight}
\end{align}
$p_s$ 和 $p_t$ 是相对于同一参考测度的两个输入密度，$w(x)$ 是密度比。支撑覆盖要求 $p_t(x)>0$ 的区域满足 $p_s(x)>0$，忽略概率为零的例外。以连续输入和离散标签为例，风险变换的关键步骤为
\begin{align}
 R_t(f)&=\int\sum_{y\in\mathcal Y}\ell(f(x),y)p_t(y\mid x)p_t(x)\,\mathrm dx\\
 &=\int\sum_{y\in\mathcal Y}\ell(f(x),y)p_s(y\mid x)
 \frac{p_t(x)}{p_s(x)}p_s(x)\,\mathrm dx.
 \label{eq:transfer-source-risk-integral}
\end{align}
第一步按输入和标签的联合概率平均损失。第二步使用协变量偏移替换条件分布，再乘除源输入密度。式\eqref{eq:transfer-source-risk-integral}因而变成带权的源域期望。需要损失可积，才能把它作为有限风险处理。源域没有出现过的输入区域没有可重用的标签信息。

给定源样本 $(x_i,y_i)_{i=1}^{n_s}$，加权经验风险为
\begin{equation}
\widehat R_t(f)=\frac1{n_s}\sum_{i=1}^{n_s}w(x_i)\ell(f(x_i),y_i).
\label{eq:transfer-weighted-empirical-risk}
\end{equation}
它增加目标常见样本对损失和梯度的贡献，不会生成新的样本。若权重估计准确，式\eqref{eq:transfer-weighted-empirical-risk}估计的是目标风险。若权重集中在少数源样本，估计方差会很大。常用诊断量是
\begin{equation}
 n_{\mathrm{eff}}=\frac{(\sum_iw_i)^2}{\sum_iw_i^2},
\end{equation}
其中 $w_i=w(x_i)$。它表示权重集中程度，不是所有损失下都精确等于独立样本数。

\subsection{密度比估计与局限}
密度比可以直接估计，也可以通过域分类器间接估计。令 $D=1$ 表示目标域，域分类器输出 $q(x)=P(D=1\mid x)$。若训练混合数据中的目标和源先验分别为 $\pi_t$ 与 $\pi_s$，贝叶斯公式给出
\begin{equation}
 w(x)=\frac{\pi_s}{\pi_t}\frac{q(x)}{1-q(x)}.
\label{eq:transfer-domain-classifier-ratio}
\end{equation}
这里 $0<\pi_t<1$且 $\pi_s=1-\pi_t$，先验是域分类器训练时的采样比例。公式还要求 $q$能可靠估计该混合分布下的后验概率。若人为把两个域采成相同数量，就必须使用相同的采样先验，而不能使用原始数据库比例。$q(x)$接近 $1$时分母很小，密度比估计会不稳定。域分类准确率很高可能意味着支撑重叠不足，不能解释为权重更可信。

一个小例子可以说明加权改变了风险。源域中低速、高速样本比例为 $(0.8,0.2)$，目标域比例为 $(0.4,0.6)$。协变量偏移条件下，两种工况的权重为 $0.5$ 和 $3$。若模型在两种工况上的平均损失为 $0.1$ 和 $0.5$，源风险为 $0.18$，加权源风险为 $0.34$，目标风险也为 $0.34$。若源域没有高速记录，则高速权重没有有限估计，重加权不能替代采集高速数据。截断大权重可以降低方差，但会引入偏差，因此必须在独立目标验证集上选择截断规则。

\subsection{类别比例变化与权重稳定性}
标签偏移时，根据式\eqref{eq:transfer-label-shift}分解联合密度可得
\begin{equation}
 p_t(x,y)=p_s(x\mid y)p_t(y)
 =p_s(x,y)\frac{p_t(y)}{p_s(y)}.
\end{equation}
因此权重应改为类别先验比 $w_y=p_t(y)/p_s(y)$，要求目标中出现的类别在源域具有正概率。这与按输入密度加权是两种不同假设。核均值匹配（Kernel Mean Matching，KMM）则通过使加权源特征均值接近目标特征均值来估计样本权重，避免分别拟合两个高维输入密度\cite{huang2007kmm}。匹配统计量仍然需要协变量偏移等条件支持其风险解释。

无标签目标数据并不直接给出真实类别比例。固定一个分类器，设共有 $K$ 类，源验证集的混淆矩阵为 $C_{ky}=P_s(\widehat Y=k\mid Y=y)$，其中 $k,y\in\{1,\ldots,K\}$。令目标预测频率为 $\mu_k=P_t(\widehat Y=k)$，目标真实类别比例为 $\pi_y^t=P_t(Y=y)$。标签偏移下，由全概率公式得到
\begin{equation}
 \mu=C\pi^t,\qquad \pi_y^t\geq0,\quad \sum_{y=1}^K\pi_y^t=1.
\end{equation}
$\mu,\pi^t\in\R^K$，$C\in\R^{K\times K}$的每列和为 $1$。可在概率单纯形约束下拟合此关系。若矩阵列近似相同，先验反演会不稳定。若类条件输入分布也改变，这个线性关系便不再成立。

例如 $C=\begin{bmatrix}0.9&0.2\\0.1&0.8\end{bmatrix}$，$\mu=(0.55,0.45)^{\mathsf T}$。写 $\pi^t=(a,1-a)^{\mathsf T}$，第一行给出 $0.55=0.9a+0.2(1-a)$，所以 $a=0.5$。预测频率 $0.55$并不等于真实类别比例 $0.5$。这一教学计算依赖混淆矩阵在目标域保持不变。

若目标中出现源域没有的类别，任何先验比都无法补充缺失类别。概念偏移下，增加样本权重也无法修复标签规则本身的变化。用真实权重计算时，对固定 $f$ 和独立源样本，式\eqref{eq:transfer-weighted-empirical-risk}在可积条件下无偏，有限方差时其方差为 $\Var_s(w\ell)/n_s$。训练后选择出的 $f$ 已依赖训练数据，不能据此宣称其训练风险是无偏泛化估计。

\section{表示统计对齐}

\subsection{MMD 的比较对象}
重要性加权需要估计密度比。高维输入中，直接估计两个密度往往困难。另一条路线是在表示空间比较源、目标分布。令编码器输出 $z=F_\theta(x)\in\mathbb R^d$，核函数 $k(u,v)$ 将表示间的相似性映射为标量。若 $k(u,v)=\langle\varphi(u),\varphi(v)\rangle_{\mathcal H}$，则最大均值差异（Maximum Mean Discrepancy，MMD）比较两个域在核特征空间中的均值：
\begin{equation}
\operatorname{MMD}^2(P_s^Z,P_t^Z)=
\left\|\mathbb E_s\varphi(Z)-\mathbb E_t\varphi(Z)\right\|_{\mathcal H}^2.
\label{eq:transfer-mmd}
\end{equation}
$\mathcal H$ 是核特征空间，$\varphi$ 是特征映射，$P_s^Z$ 与 $P_t^Z$ 是表示的边缘分布。小批次的有偏估计为
\begin{align}
\widehat{\operatorname{MMD}}^2
&=\frac1{n_s^2}\sum_{i,i'}k(z_i^s,z_{i'}^s)
+\frac1{n_t^2}\sum_{j,j'}k(z_j^t,z_{j'}^t)\\
&\quad-\frac2{n_sn_t}\sum_{i,j}k(z_i^s,z_j^t).
\label{eq:transfer-mmd-estimator}
\end{align}
源批次有 $n_s\geq1$ 个表示，目标批次有 $n_t\geq1$ 个表示。索引 $i,i'$ 遍历源批次，$j,j'$ 遍历目标批次。把式\eqref{eq:transfer-mmd}的平方范数展开为两项自身内积减去两倍交叉内积，再以样本平均替代期望，就得到式\eqref{eq:transfer-mmd-estimator}。三项分别比较源源、目标目标及源目标样本。这里包含相同下标的对角项，所以估计非负但一般有偏。

核技巧无需显式构造 $\varphi$。例如线性核 $k(u,v)=u^{\mathsf T}v$只比较原表示均值。高斯核 $k(u,v)=\exp(-\|u-v\|_2^2/(2b^2))$ 中的带宽 $b>0$控制比较尺度。采用特征核（characteristic kernel，例如适当定义域上的高斯核）且核均值存在时，总体 MMD 为零可以识别分布相同。这里的特征核具有分布可识别性，不泛指任意作用在特征上的核。有限批次 MMD 很小不证明总体分布完全相同，也不检查标签规则。

深度适配网络（Deep Adaptation Networks，DAN）将核均值差与源监督目标共同优化\cite{long2015dan}。训练时从两域各抽一个批次，计算任务损失与 MMD，再反向传播更新编码器。核带宽、批次组成和对齐权重都需要按允许的验证协议选择。

\subsection{CORAL 的二阶对齐}
均值比较未必能检查特征的波动与维度相关性。相关性对齐直接比较表示的协方差，其英文名称为 Correlation Alignment，简称 CORAL。

深度 CORAL 将这种比较作为网络的训练损失\cite{sun2016coral}。它通过比较各维波动及共同变化的统计量来调整表示。

源表示矩阵 $Z_s\in\R^{n_s\times d}$与目标表示矩阵 $Z_t\in\R^{n_t\times d}$的每行是一个样本。定义中心化矩阵 $H_s=I_{n_s}-\mathbf1\mathbf1^{\mathsf T}/n_s$，样本协方差为
\begin{equation}
 C_s=\frac{Z_s^{\mathsf T}H_sZ_s}{n_s-1},\qquad
 C_t=\frac{Z_t^{\mathsf T}H_tZ_t}{n_t-1}.
\end{equation}
$C_s,C_t\in\R^{d\times d}$是样本协方差，要求 $n_s,n_t>1$。各域样本独立同分布且具有有限二阶矩时，它们分别无偏估计该域的总体协方差。CORAL 将二者的差异定义为
\begin{equation}
 \mathcal L_{\mathrm{CORAL}}=\frac1{4d^2}\|C_s-C_t\|_F^2,
\label{eq:transfer-coral}
\end{equation}
其中 $H_t=I_{n_t}-\mathbf1\mathbf1^{\mathsf T}/n_t$与 $H_s$具有相应的批次维度。$I_n$是 $n$维单位矩阵，$\mathbf1$是相应长度的全 $1$列向量。两个 $C$均为 $d\times d$矩阵，$\|\cdot\|_F^2$表示所有矩阵元素的平方和。乘中心化矩阵 $H$相当于逐行减去批次均值。小批次协方差的秩不超过批次样本数减一，小批次会使二阶统计估计不稳定。该损失只比较二阶统计量。均值对齐需要另加损失，类别条件对齐还需要标签或可信伪标签。

例如，一维源表示为 $(-1,1)$，目标表示为 $(-2,2)$。两者均值相同，线性 MMD 为零，但方差分别为 $2$ 和 $8$，CORAL 损失为 $9$。若把源表示平移到 $(4,6)$，两者方差相同，CORAL 为零，但均值不同。这个例子说明不同统计量不能互相替代。

统计对齐的共同目标可写为
\begin{equation}
\min_{F,C}\ \mathcal L_s\bigl(C(F(X_s)),Y_s\bigr)
+\lambda\,\mathcal L_{\mathrm{align}}\bigl(F(X_s),F(X_t)\bigr),
\label{eq:transfer-statistical-alignment}
\end{equation}
这里 $X_s,Y_s$表示带标签源批次，$X_t$表示无标签目标批次，$C$是任务预测头，$\mathcal L_s$是源批次平均监督损失。$\mathcal L_{\mathrm{align}}$可以取上述 MMD 或 CORAL，$\lambda\geq0$控制对齐强度。采用 MMD 的目标沿用标签\eqref{eq:transfer-mmd-objective}。
\begin{equation}
 \mathcal L=\mathcal L_s+\lambda\widehat{\operatorname{MMD}}^2.
 \label{eq:transfer-mmd-objective}
\end{equation}
$\lambda$过大可能把不同类别的表示拉近。边缘分布对齐下降不能证明 $P_s(Y\mid Z)=P_t(Y\mid Z)$。图\ref{fig:transfer-statistics}用一维点集区分均值和方差，图中没有进行逐个样本配对。

\begin{figure}[htbp]
\centering
\begin{tikzpicture}[x=0.62cm,y=0.8cm,font=\small]
\draw[->] (-3,0)--(7,0) node[right] {$z$};
\foreach \v in {-2,-1,0,1,2,4,5,6}{\draw (\v,0.06)--(\v,-0.06) node[below,font=\scriptsize] {$\v$};}
\node[anchor=east] at (-3,1) {源};
\node[anchor=east] at (-3,2) {目标 A};
\node[anchor=east] at (-3,3) {目标 B};
\draw[dashed,gray] (0,0.3)--(0,3.35);
\foreach \v in {-1,1}{\fill (\v,1) circle (2.5pt);}
\foreach \v in {-2,2}{\draw[thick] (\v-0.09,1.9) rectangle (\v+0.09,2.1);}
\foreach \v in {4,6}{\draw[thick] (\v,3) circle (2.5pt);}
\node[anchor=west,font=\scriptsize] at (2.6,1) {均值 $0$，方差 $2$};
\node[anchor=west,font=\scriptsize] at (2.6,2) {均值 $0$，方差 $8$};
\node[anchor=west,font=\scriptsize] at (-2.5,3.7) {目标 B 均值 $5$，方差 $2$};
\end{tikzpicture}
\caption{统计量比较的教学点集。实心点、方框和空心圆分别表示源、目标 A 和目标 B。方差按两条样本的无偏样本方差计算。源与 A 的均值相同，源与 B 的方差相同，两种匹配均不能识别所有分布差异。}
\label{fig:transfer-statistics}
\end{figure}

\section{域对抗适配}

\subsection{DANN 的信息流}
统计对齐预先规定了比较哪些统计量。域对抗神经网络（Domain-Adversarial Neural Network，DANN）改用一个可训练的域判别器 $D$寻找可区分源、目标的表示差异，同时要求编码器保持源域预测能力\cite{ganin2016dann}。用 $a=0$表示源域、$a=1$表示目标域，以避免与表示维度 $d$混淆。判别器输出 $D(z)\in(0,1)$表示目标域概率。采用两域等权的交叉熵损失
\begin{equation}
 \mathcal L_d=-\mathbb E_s\log(1-D(Z))-\mathbb E_t\log D(Z).
 \label{eq:transfer-domain-loss}
\end{equation}
这里两项分别惩罚源样本被判为目标、目标样本被判为源。实际用两域各自的批次平均近似。理想的极小极大目标为
\begin{equation}
\min_{F,C}\max_D\quad\mathcal L_s-\lambda\mathcal L_d,
\qquad\lambda\geq0.
\label{eq:transfer-dann}
\end{equation}
域判别器最大化负的域损失，等价于最小化自身交叉熵。编码器最小化整个目标，因此会使域判别变难。实际实现常在编码器与域判别器之间放置梯度反转层。图\ref{fig:transfer-dann}说明两条梯度路径的方向差异。

\begin{figure}[htbp]
\centering
\begin{tikzpicture}[node distance=0.7cm and 0.8cm]
\node[idi input] (x) {输入 $x$};
\node[idi process,right=of x] (f) {编码器\\$z=F_\theta(x)$};
\node[idi output,right=of f] (c) {任务头 $C$\\源标签损失};
\node[idi output,below=of c] (d) {域判别器 $D$\\域损失};
\node[idi process,below=of f] (r) {梯度反转\\前向不变，反向乘 $-\lambda$};
\draw[idi arrow] (x)--(f)--(c);
\draw[idi arrow] (f)--(r)--(d);
\draw[idi feedback] (d.south) -- ++(0,-0.4) -| (r.south);
\end{tikzpicture}
\caption{DANN 的两条训练路径。任务路径要求源域可预测，域路径通过梯度反转要求表示减少域信息。}
\label{fig:transfer-dann}
\end{figure}

若编码器参数为 $\theta$，判别器参数为 $\omega$，任务头参数为 $\psi$，梯度下降实际执行
\begin{align}
 \theta&\leftarrow\theta-\eta(\nabla_\theta\mathcal L_s-\lambda\nabla_\theta\mathcal L_d),\\
 \omega&\leftarrow\omega-\eta_d\nabla_\omega\mathcal L_d,
 \qquad\psi\leftarrow\psi-\eta_c\nabla_\psi\mathcal L_s.
\end{align}
$\eta,\eta_d,\eta_c>0$是各模块学习率。编码器与判别器对域损失使用相反方向。源标签监督用于防止所有表示都变成一个常量，但它仍不能单独保证目标分类正确。

固定编码器，且表示具有相对于同一参考测度的密度时，式\eqref{eq:transfer-domain-loss}在每个 $z$处对 $D$的导数为
\begin{equation}
 \frac{p_s(z)}{1-D(z)}-\frac{p_t(z)}{D(z)}.
\end{equation}
将其设为零，得到等先验下的理想判别器
\begin{equation}
 D^*(z)=\frac{p_t(z)}{p_s(z)+p_t(z)}.
\end{equation}
此式要求分母大于零，只有一域具有密度时按边界极限理解。足够强且充分训练的判别器接近随机，才支持边缘分布难以区分的解释。欠拟合判别器也可能接近随机。即使边缘分布完全相同，也没有检验标签关系。

例如两域的 $Z$均匀取 $-1,+1$。源域标签规则为 $Y=\mathbf1[Z=1]$，目标域规则为 $Y=\mathbf1[Z=-1]$，$\mathbf1[\cdot]$为条件成立时取 $1$的指示函数。域判别器只能猜测，而按源规则分类的模型会在目标域全部判错。

领域适配理论用目标风险的上界表达这一限制\cite{bendavid2010theory}。对二分类零一损失，令 $\mathcal H$为分类器集合，$h\in\mathcal H$，则总体风险满足
\begin{equation}
R_t(h)\le R_s(h)+\frac12d_{\mathcal H\Delta\mathcal H}(P_s^X,P_t^X)+\lambda^*,
\label{eq:transfer-domain-bound}
\end{equation}
这里 $R_s,R_t$是源、目标总体错误率，$P_s^X,P_t^X$是输入边缘分布。约定差异量为
\begin{equation}
 d_{\mathcal H\Delta\mathcal H}(P_s^X,P_t^X)
 =2\sup_{h,h'\in\mathcal H}
 \left|P_s(h(X)\ne h'(X))-P_t(h(X)\ne h'(X))\right|.
\end{equation}
它比较两个分类器在两域上意见不一致的概率，$\sup$取所有分类器对中的上确界。$\lambda^*$是同一个分类器在两域的最小风险之和，若最小值不取得则以相应下确界定义。它不同于训练超参数 $\lambda$。式\eqref{eq:transfer-domain-bound}是总体风险界，若改用有限样本估计风险和差异，还需加入相应的估计与复杂度项。DANN 的域损失不等于此差异量的精确值。对齐降低差异项时，不能保证 $\lambda^*$很小。

如果目标类别比例改变，盲目匹配边缘表示还可能破坏类别结构。条件对齐将类别信息纳入比较，试图区分“同类跨域差异”与“不同类差异”。目标伪标签不可靠时，这种比较会传播错误。若源域包含目标没有的类别，需要降低源独有类的影响。若目标出现源域没有的类别，则需要未知类识别或新标签证据。上述变化修改了标签空间条件，不能仅靠增大对齐系数解决。

\section{参数适配与任务变化}

\subsection{冻结、微调与适配器}
前面的重加权和对齐主要依赖同一任务的源标签。如果目标任务改变，最直接的证据是目标标签。任务适配首先要决定哪些参数可以改变。沿用第\ref{chap:pretraining}章的记法，编码器为 $f_\theta$，目标头为 $h_\psi$。它们分别对应前文的 $F$和 $C$。给定 $m\geq1$个带标签目标训练样本 $(x_i^t,y_i^t)$，监督适配目标为
\begin{equation}
 \mathcal L_t(\theta,\psi)=\frac1m\sum_{i=1}^m
 \ell\bigl(h_\psi(f_\theta(x_i^t)),y_i^t\bigr).
\end{equation}
$\theta$从预训练值 $\theta_0$出发，$\psi$是目标头参数。分类时头输出类别概率，$\ell$可取交叉熵。回归时输出和损失应按目标量定义。若头输出概率，预测类别还需经最大概率等决策规则得到，不能把概率向量和类别编号混为同一对象。

线性探测（linear probing）固定 $\theta=\theta_0$，只训练对表示为仿射映射的类别得分头。它检验表示中的线性可分信息。全量微调（fine-tuning）同时更新 $\theta$和 $\psi$，自由度最大，但目标标签少时容易过拟合和破坏原有表示。逐层解冻先训练高层，再用验证数据判断是否需要改变较低层。冻结还应说明是否固定归一化运行统计，否则权重没有梯度也可能改变模型函数。比较这些策略时，应固定目标标签量、预训练模型、数据划分和验证预算。

Adapter 在冻结骨干中加入小型瓶颈模块。对 $h\in\mathbb R^d$，一个残差 Adapter 为
\begin{equation}
 h'=h+W_{\mathrm{up}}\sigma(W_{\mathrm{down}}h),
\end{equation}
其中 $W_{\mathrm{down}}\in\mathbb R^{r\times d}$、$W_{\mathrm{up}}\in\mathbb R^{d\times r}$，$1\leq r<d$，$\sigma$是逐元素非线性函数。$h'$与 $h$均为 $d$维，残差连接保留原计算路径。忽略偏置时新增 $2dr$个参数。Adapter 将目标变化限制在小型非线性模块中\cite{houlsby2019adapter}。只训练新增模块和必要的目标头可以为不同目标域保存不同版本。若瓶颈过窄，所需变化可能无法表示。模块数量增加后需要路由规则，非线性模块通常不能简单合并回骨干。

\subsection{LoRA 的低秩更新}
Adapter 新增非线性计算路径。若希望把目标变化直接写为线性权重增量，可以使用低秩适配（Low-Rank Adaptation，LoRA）\cite{hu2021lora}。它冻结原权重 $W_0\in\mathbb R^{d_{\rm out}\times d_{\rm in}}$，用低秩矩阵表示更新：
\begin{equation}
 y=(W_0+sBA)x,
\qquad A\in\mathbb R^{r\times d_{\rm in}},\quad
 B\in\mathbb R^{d_{\rm out}\times r},\quad s=\alpha/r.
\label{eq:transfer-lora}
\end{equation}
$x\in\mathbb R^{d_{\rm in}}$ 是输入，$y\in\mathbb R^{d_{\rm out}}$ 是输出，$r$ 是更新秩，$\alpha$ 是尺度超参数。新增参数数目为 $r(d_{\rm in}+d_{\rm out})$，更新矩阵的秩不超过 $r$。实际节省参数要求该数量小于 $d_{\rm in}d_{\rm out}$。低秩假设约束的是有用的权重变化，不是说原权重本身低秩。

若 $G=\partial\mathcal L/\partial W_{\rm eff}$ 且 $W_{\rm eff}=W_0+sBA$，链式法则给出
\begin{equation}
\nabla_B\mathcal L=sGA^{\mathsf T},\qquad
\nabla_A\mathcal L=sB^{\mathsf T}G.
\end{equation}
$G$与 $W_0$同形，两个梯度分别与 $B$和 $A$同形。常用初始化令 $A$随机、$B=0$，使初始更新为零但 $B$可以收到非零梯度。若两者均为零，两条梯度都为零，低秩分支无法启动。

用 $d_{\rm in}=d_{\rm out}=2$、$r=1$、$s=1$的教学矩阵核对运算。取 $A=[1\ \ 2]$、$B=[3\ \ 4]^{\mathsf T}$，则 $BA=\begin{bmatrix}3&6\\4&8\end{bmatrix}$。两列成比例，更新秩为 $1$。输入 $x=(1,0)^{\mathsf T}$时，新增输出为 $(3,4)^{\mathsf T}$。这个极小例中仍有 $4$个新增参数，没有节省。原层为 $100\times100$且 $r=2$时，才由 $10000$个全量权重更新变为 $400$个低秩参数。

图\ref{fig:transfer-lora}把低秩更新与原层分开。部署时若目标版本固定，可将 $sBA$合并进 $W_0$，不再显式计算低秩分支。多个目标版本需要保留或切换各自增量。冻结骨干仍需前向计算，反向传播也可能需要保存中间激活。可训练参数、峰值训练显存和推理时延因此应分别报告。
\begin{figure}[htbp]
\centering
\begin{tikzpicture}[node distance=0.6cm and 0.6cm]
\node[idi input] (x) {$x\in\R^{d_{\rm in}}$};
\node[idi state,right=of x] (w) {$W_0$\\冻结};
\node[idi output,right=1.2cm of w] (sum) {相加\\输出 $y$};
\node[idi process,below=of w] (a) {$A$ 可训练\\降至 $r$维};
\node[idi process,right=of a] (b) {$sB$ 可训练\\升至 $d_{\rm out}$维};
\draw[idi arrow] (x)--(w)--(sum);
\draw[idi arrow] (x.south) |- (a.west);
\draw[idi arrow] (a)--(b);
\draw[idi arrow] (b.north)--(sum.south);
\end{tikzpicture}
\caption{LoRA 的局部计算。实线箭头表示前向计算，冻结路径输出 $W_0x$，可训练路径输出 $sBAx$，相加得到式\eqref{eq:transfer-lora}。图中的 $r$是中间宽度，不是原权重的秩。}
\label{fig:transfer-lora}
\end{figure}

LoRA 降低存储和优化参数量，但小秩可能表达不了目标任务所需的变化。参数高效不等于风险自动更低，也不等于持续学习中的旧能力保持。连续覆盖同一个 Adapter 或 LoRA 仍会遗忘。为不同任务保存模块则还需正确路由，不能仅凭模块冻结宣称整个系统的旧预测不变。

\section{测试时适配与负迁移}

\subsection{无标签目标数据的在线使用}
测试时适配（test-time adaptation，TTA）在部署输入到达后更新有限参数。它必须规定适配发生在预测前还是预测后，以及更新状态是否跨批次保留。以熵最小化为例，模型输出 $p_\eta(k\mid x)$，允许更新的参数为 $\eta$，批次 $B$ 上的目标为
\begin{equation}
\min_\eta\frac1{|B|}\sum_{x\in B}
\left[-\sum_{k=1}^{K}p_\eta(k\mid x)\log p_\eta(k\mid x)\right].
\label{eq:transfer-tta-entropy}
\end{equation}
$B$是非空部署批次，$|B|$是样本数，$K\geq2$是固定类别数，$p_\eta(k\mid x)$是第 $k$类预测概率，约定 $0\log0=0$，对数取自然对数。Tent 在源训练完成后，结合批次归一化统计，按其设置更新归一化层的仿射参数来降低预测熵\cite{wang2021tent}。它可以不访问源训练样本，但仍使用源训练模型以及关于目标预测结构的假设。

二分类预测 $(0.5,0.5)$的熵约为 $0.693$，$(0.9,0.1)$的熵约为 $0.325$。若真实标签是第二类，后一个预测更自信但仍错误。进一步令第一类概率为 $p=\sigma(a)=1/(1+e^{-a})$，$a\in\R$为二分类得分差，熵为 $H=-p\log p-(1-p)\log(1-p)$。由链式法则
\begin{equation}
 \frac{\mathrm dH}{\mathrm da}
 =\log\frac{1-p}{p}\,p(1-p).
\end{equation}
当 $p>1/2$时导数为负，梯度下降继续增大 $a$，使第一类概率更高。没有标签时，这个目标没有判断当前偏向是否正确的依据。熵最小化因而依赖初始预测大体可信等条件。类别不平衡、未知类别或小批次下，它可能强化错误自信或使输出集中到少数类。

严格在线协议应先用当前参数预测并记录结果，再用当前批次无标签输入更新后续预测。若先适配再预测，则允许本批输入影响其自身预测，需要单独报告这种协议。每批重置到源参数与跨批累积更新也不相同。两种方式都不能访问未来批次和当前测试标签。步数、学习率、重置频率及停止规则应由独立开发数据或预先规定的规则确定。无标签熵下降不能作为最终验收标准。

\subsection{负迁移的判定}
负迁移指使用源知识后，目标风险高于合适的无迁移基线。设迁移方法和基线得到的预测器分别为 $f_{\rm tr}$与 $f_{\rm base}$，定义
\begin{equation}
 \Delta R_t=R_t(f_{\rm tr})-R_t(f_{\rm base}).
\end{equation}
$\Delta R_t>0$表示按当前损失发生负迁移。基线必须有相同的目标标签权限、数据划分和合理的优化预算。有限测试集上的微小差值还可能来自抽样与训练随机性，需要报告不确定性，不能直接当作确定的总体结论。

负迁移可能来自标签含义变化、表示丢失目标信息、错误的域对齐或过强的参数约束。它与遗忘的评价对象不同。负迁移检查源知识是否损害目标学习。遗忘检查更新以后旧任务是否退化。目标改善和源域遗忘可以同时发生，目标负迁移也可以出现在源预测完全没有改变的冻结表示方法中。

\section{迁移评价与方法选择}

\subsection{信息权限与验证划分}
评价首先固定数据权限。普通监督任务适配使用目标训练标签选择参数，再用独立验证集选择策略。UDA 使用源标签与目标无标签训练记录，若另外使用目标验证标签，必须说明其数量和用途。严格不允许目标标签的设置，应在独立开发域上确定配置或采用预先规定的无标签选择规则。不能训练时声称无目标标签，选择最优检查点时却查看目标测试准确率。

DG 应按整个域留出目标。用其无标签记录计算归一化、密度比或对齐损失，都会改变信息条件。TTA 则明确允许当前部署输入参与更新，需按时间顺序记录成绩。这些试验都可研究迁移，却不能将不同数据权限的分数当作同条件算法比较。

图\ref{fig:transfer-evaluation}把适配和选择的边界分开。对振动信号先按设备、故障事件或连续时间段划分，再在各自范围内切窗口。同一事件的重叠窗口不能分到训练与测试两侧。若是跨相机试验，也应避免同一产品的近重复图像跨集合。
\begin{figure}[htbp]
\centering
\begin{tikzpicture}[node distance=0.7cm and 0.6cm]
\node[idi input,text width=2.1cm] (train) {允许的训练数据\\拟合参数};
\node[idi input,text width=2.1cm,right=of train] (val) {开发域或验证集\\选择配置};
\node[idi output,text width=2.1cm,right=of val] (test) {冻结选择后\\独立目标评价};
\draw[idi arrow] (train)--(val)--(test);
\node[below=0.5cm of val,text width=7.5cm,align=center,font=\small] {无监督适配可访问声明允许的目标输入\\目标测试标签不反向进入参数拟合或配置选择};
\end{tikzpicture}
\caption{迁移评价的信息边界。箭头表示训练、选择、测试的顺序，不表示允许把测试标签传回训练。是否使用目标测试输入，须按 UDA 或 TTA 的具体协议另行声明。}
\label{fig:transfer-evaluation}
\end{figure}

\subsection{基线、指标与资源}
有目标训练标签时，至少比较从零训练、线性探测、全量微调和适用的轻量适配。在多个目标标签预算下绘制学习曲线，可以检查源知识是否主要帮助小样本情况。无目标标签时，从零监督训练目标模型不构成可行基线，应先比较源域直接训练后不适配的模型。若使用目标标签训练作为参考，应明确标为额外监督条件。

评价不能只报告对齐损失、域分类准确率或训练损失。应在固定目标测试集上报告所选损失的经验风险，以及与业务有关的各类召回和误报。概率输出还应检查校准，即相同置信度预测的实际正确比例是否与置信度相符。标签比例失衡时不能只报告总准确率。DG 应报告各个留出域而不仅是合并平均分。若包含未知类别，拒识阈值应在验证阶段确定，并分别报告已知类识别与未知类检出。

重复运行时保持数据划分一致，报告均值及波动范围。对于同一设备内高度相关的窗口，应按独立设备或事件组织重采样和区间估计。训练时间、可训练参数、总参数、峰值显存和推理时延分别计量。共享预训练的成本可以摊销，但仍要说明来源和是否计入比较。

\subsection{新设备的适配流程}
回到新相机质检。先核对新旧缺陷定义是否一致，并记录亮度、分辨率和产品类别的变化。有少量目标标签时，先建立线性探测与小学习率微调基线。若预算只允许保存少量目标专用参数，再比较 Adapter 或 LoRA。微调策略的选择仍由验证成绩决定。

若只有目标无标签图像，首先保留不适配的源模型作为对照。只有在标签空间一致且可共享预测关系的假设合理时，才进一步比较密度比加权或表示对齐。输入重叠很差时应补充数据，不依靠极端权重解决覆盖缺口。缺陷定义发生变化或出现新故障时，应获得相应标签，不能把边缘对齐当成新标签证据。TTA 还要预定重置与回退规则，防止连续错误更新。

更新后若系统仍服务旧相机，应同时报告旧相机性能变化。目标性能上升而源域下降时，是否接受这种取舍属于预先规定的服务目标。第\ref{chap:continual}章将把这种新旧能力冲突扩展到连续到达的数据流，并用逐阶段评价检查学习和遗忘。

\section*{辅助学习材料}
\begingroup
\emergencystretch=3em
\begin{itemize}
\item \textbf{Transfer Learning}。作者/平台：Stanford CS231n。英文课程笔记。阅读固定特征提取和微调的选择条件，配合本章的参数适配部分使用。\href{https://cs231n.github.io/transfer-learning/}{在线阅读}。
\item \textbf{Transfer Learning for Computer Vision Tutorial}。作者/平台：PyTorch。英文代码教程。比较冻结特征提取器和全量微调流程。\href{https://docs.pytorch.org/tutorials/beginner/transfer\_learning\_tutorial.html}{在线阅读}。
\item \textbf{什么是迁移学习（Transfer Learning）？}。知乎中文解说。对照冻结特征与微调的可训练参数范围。标题与主题沿用项目记录，原页受限。\href{https://zhuanlan.zhihu.com/p/1899236885408256815}{在线阅读}。
\item \textbf{迁移学习(Transfer Learning)概述及代码实现}。CSDN 中文博客。对照源域与目标域划分，用无迁移基线检查负迁移。标题与主题沿用项目记录，原页受限。\href{https://blog.csdn.net/weixin\_49272172/article/details/115791703}{在线阅读}。
\item \textbf{第十一节 2021 - 概述领域自适应 (Domain Adaptation)}。李宏毅中文课程，第110集。对照分布差异与适配条件。沿用项目课程链接，本次未逐帧核验视频。\href{https://www.bilibili.com/video/BV1Wv411h7kN/?p=110}{观看分集}。
\end{itemize}
\endgroup

\paragraph{本章小结}
迁移学习先要确定改变的是任务、域还是两者。协变量偏移允许用输入密度比重加权。统计对齐和域对抗改变表示分布，但不能独立证明标签关系保持。冻结、微调、Adapter 和 LoRA 改变可训练参数的范围，容量越小并不自动意味着目标风险越低。测试时适配可以利用无标签输入，却需要明确预测与更新顺序。所有方法都应与无迁移基线和固定目标测试集比较，并检查负迁移。若目标数据继续到达且旧能力仍需保留，下一章的持续学习需要在迁移评价之外记录逐阶段的学习和遗忘。
